\documentclass[a4paper,12pt]{article}
\usepackage{iftex}
\ifPDFTeX
  \usepackage[T2A]{fontenc}
  \usepackage[utf8]{inputenc}
  \usepackage[russian]{babel}
  \usepackage{helvet}
  \renewcommand{\familydefault}{\sfdefault}
\else
  \usepackage{fontspec}
  \setmainfont{Liberation Sans}
  \usepackage[russian]{babel}
\fi
\usepackage{amsmath,amssymb,mathtools}
\usepackage{geometry}
\geometry{left=22mm,right=22mm,top=21mm,bottom=24mm}
\usepackage{microtype}
\usepackage{booktabs,tabularx,array}
\usepackage{enumitem}
\usepackage{tikz}
\usepackage{xcolor}
\usepackage{fancyhdr}
\usepackage{hyperref}
\usepackage{parskip}
\usepackage{needspace}
\definecolor{accent}{HTML}{176E78}
\definecolor{dark}{HTML}{18383E}
\definecolor{muted}{HTML}{56666B}
\definecolor{err}{HTML}{A12D39}
\definecolor{good}{HTML}{157460}
\hypersetup{colorlinks=true,linkcolor=accent,urlcolor=accent,pdfauthor={Лёвин Артём Александрович},pdftitle={Ревью домашней работы — Анастасия Павлова — 29 октября 2026 г.}}
\setlength{\parindent}{0pt}
\setlength{\parskip}{5.5pt}
\renewcommand{\arraystretch}{1.25}
\setlist[itemize]{leftmargin=*,itemsep=2pt,topsep=3pt}
\pagestyle{fancy}
\fancyhf{}
\fancyfoot[C]{\footnotesize\color{muted}Лёвин Артём Александрович эксклюзивно для Анастасии Павловой}
\fancyfoot[R]{\footnotesize\color{muted}\thepage}
\renewcommand{\headrulewidth}{0pt}
\renewcommand{\footrulewidth}{0pt}
\newcommand{\blocktitle}[1]{\needspace{4\baselineskip}\vspace{8pt}{\large\bfseries\color{dark}#1}\par\vspace{2pt}}
\newcommand{\errorlabel}{\textcolor{err}{\textbf{Ошибка:}}}
\newcommand{\keylabel}{\textcolor{accent}{\textbf{Ключевая идея:}}}
\newcommand{\resultlabel}{\textcolor{good}{\textbf{Итог:}}}
\begin{document}
\begin{titlepage}
\thispagestyle{empty}
\centering
\vspace*{36mm}
{\small\bfseries\color{accent}МАТЕМАТИКА\par}
\vspace{13mm}
{\Huge\bfseries\color{dark}Ревью\par}
\vspace{2mm}
{\LARGE\bfseries\color{dark}домашней работы\par}
\vspace{9mm}
{\large Кредиты с равномерным погашением основного долга\par}
\vspace{23mm}
{\large 29 октября 2026 г.\par}
\vspace{18mm}
{\normalsize Автор: Лёвин Артём Александрович\par}
\vfill
{\small\color{muted}Лёвин Артём Александрович\par}
{\small\color{muted}эксклюзивно для Анастасии Павловой\par}
\vspace{8mm}
\begin{tikzpicture}[x=1cm,y=1cm]
\draw[accent!55,line width=0.85pt] (0,0) .. controls (2.3,0.6) and (5.2,-0.65) .. (9.7,0.15);
\end{tikzpicture}
\end{titlepage}
\setcounter{page}{2}
\section*{Сводная таблица}
\vspace{3mm}
В таблице собраны задания, в которых требовалось дополнить решение или исправить итоговую запись. Номер задания ведёт к подробному разбору.
\par\vspace{6mm}\noindent
{\small
\setlength{\tabcolsep}{4pt}
\begin{tabularx}{\linewidth}{@{}>{\centering\arraybackslash}p{9mm}>{\raggedright\arraybackslash}p{26mm}>{\raggedright\arraybackslash}X>{\raggedright\arraybackslash}p{27mm}>{\raggedright\arraybackslash}p{35mm}@{}}
\toprule
\textbf{№} & \textbf{Тема} & \textbf{Суть ошибки} & \textbf{Ваш ответ} & \textbf{Правильный ответ}\\
\midrule
\hyperref[task:2]{\textbf{2}} & Сумма ежегодных платежей & Платежи по годам найдены; итоговая сумма не вычислена. & \(78;\ 72;\ 66\) тыс. руб. & \(78;\ 72;\ 66\) тыс. руб.; всего \(216\) тыс. руб.\\
\addlinespace[5pt]
\hyperref[task:9]{\textbf{9}} & Диапазон процентных ставок & Границы для доли \(p\) найдены верно, но ошибочно записаны проценты. & \(10\%\); верхняя граница записана с ошибкой в знаке \(\%\). & От \(10\%\) до \(11\frac23\%\) включительно.\\
\bottomrule
\end{tabularx}
}
\vfill
{\small\color{muted}Рассмотрены только задания, требующие исправления или дополнения.}
\clearpage
\thispagestyle{empty}
\vspace*{\fill}
\begin{center}
{\Huge\bfseries\color{dark}Разбор ошибок}
\end{center}
\vspace*{\fill}
\clearpage
\phantomsection\label{task:2}
\section*{Задание № 2. Общая сумма платежей}
\blocktitle{Текст задания}
Кредит \(180\) тыс. руб. выдан на \(3\) года под \(10\%\) годовых. Основной долг погашается равными долями. Составьте таблицу из пяти колонок, найдите ежегодные платежи и общую сумму платежей.

\blocktitle{Ваше решение}
Вы верно указали значения для каждого года. По записи в тетради получаем следующую таблицу (все суммы в тыс. руб.):
\begin{center}
\small
\setlength{\tabcolsep}{5pt}
\begin{tabular}{@{}ccccc@{}}
\toprule
\textbf{Год} & \textbf{Долг в начале} & \textbf{Проценты} & \textbf{Платёж} & \textbf{Остаток}\\
\midrule
1 & 180 & 18 & 78 & 120\\
2 & 120 & 12 & 72 & 60\\
3 & 60 & 6 & 66 & 0\\
\bottomrule
\end{tabular}
\end{center}

\textbf{Ваш ответ:} ежегодные платежи \(78\), \(72\), \(66\) тыс. руб.; общая сумма выплат отдельно не указана.

\blocktitle{Где именно возникает недочёт}
\textbf{Последний правильный шаг:} третий платёж \(66\) тыс. руб., остаток долга после него равен нулю. Таблица составлена верно.

\textbf{Первый неполный шаг:} после таблицы нужно было сложить все ежегодные платежи. Это требование содержится в условии, но вычисления итоговой суммы нет.

\errorlabel\ Отсутствует обязательный результат: общая сумма всех выплат за три года.

\blocktitle{Почему важно выполнить последний шаг}
Ежегодный платёж состоит из двух частей: погашения части основного долга и процентов за пользование оставшейся суммой. Поэтому сумма всех платежей больше первоначального кредита. В данном случае одних значений по годам недостаточно для полного ответа.

\keylabel\ Если известны все платежи, общая выплаченная сумма равна их сумме. Проверить результат можно также как сумму первоначального долга и всех начисленных процентов.

\blocktitle{Исправление от последнего правильного шага}
Таблицу менять не требуется. Сложим её значения в столбце «Платёж»:
\[
78+72+66=\boxed{216\ \text{тыс. руб.}}
\]
\clearpage
\section*{Задание № 2. Подробное решение (продолжение)}
\blocktitle{Полное правильное решение}
\textbf{Шаг 1.} Основной долг каждый год уменьшается на одну и ту же сумму:
\[
180:3=60\ \text{тыс. руб.}
\]
\textbf{Шаг 2.} Проценты за первый год начисляются на \(180\) тыс. руб.:
\[
180\cdot 0{,}10=18,\qquad P_1=60+18=78.
\]
После платежа осталось \(180-60=120\) тыс. руб. основного долга.

\textbf{Шаг 3.} Во втором году проценты начисляются уже на \(120\) тыс. руб.:
\[
120\cdot 0{,}10=12,\qquad P_2=60+12=72.
\]
Остаток основного долга равен \(120-60=60\) тыс. руб.

\textbf{Шаг 4.} В третьем году:
\[
60\cdot 0{,}10=6,\qquad P_3=60+6=66.
\]
После последнего платежа основной долг погашен полностью.

\textbf{Шаг 5.} Найдём общую сумму выплат:
\[
P_1+P_2+P_3=78+72+66=\boxed{216\ \text{тыс. руб.}}
\]

\blocktitle{Проверка другим способом}
Сумма погашений основного долга равна \(3\cdot60=180\) тыс. руб. Сумма процентов:
\[
18+12+6=36\ \text{тыс. руб.}
\]
Тогда вся выплаченная сумма равна \(180+36=216\) тыс. руб. Оба способа дали одинаковый результат.

\resultlabel\ Ежегодные платежи: \(78\), \(72\), \(66\) тыс. руб.; общая сумма выплат: \(\boxed{216\ \text{тыс. руб.}}\).

\blocktitle{Задача на закрепление}
Кредит \(240\) тыс. руб. выдан на \(4\) года под \(5\%\) годовых с равномерным погашением основного долга. Составьте таблицу из пяти колонок: год, долг на начало года, начисленные проценты, ежегодный платёж и остаток основного долга. Найдите все ежегодные платежи и \textbf{общую сумму выплат}.
\clearpage
\phantomsection\label{task:9}
\section*{Задание № 9. Границы годовой ставки}
\blocktitle{Текст задания}
Кредит \(3\) млн руб. выдан на \(6\) лет с равномерным погашением основного долга. Наибольший платёж не превышает \(0{,}85\) млн руб., а наименьший составляет не менее \(0{,}55\) млн руб. Найдите все возможные годовые процентные ставки.

\blocktitle{Ваше решение}
Вы правильно составили неравенства для первого и последнего платежей и получили:
\[
3p\le 0{,}35,\qquad 0{,}5p\ge 0{,}05;
\]
\[
\boxed{0{,}1\le p\le\frac7{60}}.
\]
При этом в заключительной записи верхняя граница представлена дробью с символом процента; из этой записи получается другое числовое значение.

\textbf{Ваш ответ:} нижняя граница \(10\%\) указана верно; верхняя граница в процентах записана неверно.

\blocktitle{Где именно возникает ошибка}
\textbf{Последний правильный шаг:} \(0{,}1\le p\le\frac7{60}\), где \(p\) означает \emph{долю единицы}, используемую в формулах для процентов.

\textbf{Первый неправильный шаг:} перевод \(\frac7{60}\) в запись со знаком \(\%\). Нельзя просто приписать символ процента к десятичной дроби или обыкновенной дроби, если она уже обозначает долю единицы.

\errorlabel\ Смешаны два способа записи ставки: доля единицы и проценты. Например, \(\frac7{60}\) как доля единицы не равно \(\frac7{60}\%\).

\blocktitle{Почему это неверно}
Знак \(\%\) означает «одна сотая». Поэтому \(1\%=0{,}01\). Чтобы выразить долю \(p\) в процентах, нужно умножить \(p\) на \(100\) и затем записать знак \(\%\):
\[
\frac7{60}\cdot100\%=\frac{700}{60}\%=\boxed{11\frac23\%}.
\]
Если написать \(\frac7{60}\%\), получится лишь \(\frac7{6000}\) доли единицы, что в \(100\) раз меньше нужной верхней границы.
\clearpage
\section*{Задание № 9. Подробное решение (продолжение)}
\keylabel\ Сначала найдите диапазон для \(p\) как доли единицы. Только после этого переведите \emph{обе} границы в проценты умножением на \(100\).

\blocktitle{Полное правильное решение}
\textbf{Шаг 1.} Ежегодная часть погашаемого основного долга:
\[
\frac{3}{6}=0{,}5\ \text{млн руб.}
\]
\textbf{Шаг 2.} В первом году проценты начисляются на весь кредит \(3\) млн руб. Поэтому наибольший платёж:
\[
P_1=0{,}5+3p.
\]
В начале шестого года остаётся \(0{,}5\) млн руб. основного долга. Поэтому наименьший платёж:
\[
P_6=0{,}5+0{,}5p.
\]
\textbf{Шаг 3.} Переведём условия на платежи в неравенства:
\[
\begin{cases}
0{,}5+3p\le0{,}85,\\
0{,}5+0{,}5p\ge0{,}55.
\end{cases}
\]
Вычтем \(0{,}5\) из обеих частей каждого неравенства:
\[
\begin{cases}
3p\le0{,}35,\\
0{,}5p\ge0{,}05.
\end{cases}
\]
Разделим на положительные числа \(3\) и \(0{,}5\); знаки неравенств сохраняются:
\[
\begin{cases}
p\le\dfrac{0{,}35}{3}=\dfrac7{60},\\[3pt]
p\ge\dfrac{0{,}05}{0{,}5}=0{,}1=\dfrac1{10}.
\end{cases}
\]
\textbf{Шаг 4.} Переведём найденный промежуток в проценты:
\[
\begin{aligned}
0{,}1\cdot100\%&=10\%,\\
\frac7{60}\cdot100\%&=11\frac23\%.
\end{aligned}
\]
Знаки \(\le\) и \(\ge\) допускают равенство, поэтому обе границы включаются.
\clearpage
\section*{Задание № 9. Проверка и закрепление}
\blocktitle{Проверка найденных границ}
\textbf{Нижняя граница:} \(p=0{,}1\) (то есть \(10\%\)). Получаем:
\[
P_1=0{,}5+3\cdot0{,}1=0{,}8\le0{,}85,
\]
\[
P_6=0{,}5+0{,}5\cdot0{,}1=0{,}55\ge0{,}55.
\]
Оба условия выполнены.

\textbf{Верхняя граница:} \(p=\frac7{60}\) (то есть \(11\frac23\%\)). Тогда:
\[
P_1=0{,}5+3\cdot\frac7{60}=0{,}85\le0{,}85,
\]
\[
P_6=0{,}5+0{,}5\cdot\frac7{60}=\frac{67}{120}\approx0{,}5583\ge0{,}55.
\]
Оба условия снова выполнены. Любая промежуточная ставка также подходит, поскольку она удовлетворяет обоим полученным неравенствам.

\resultlabel\ Все возможные годовые ставки:
\[
\boxed{10\%\le r\le 11\frac23\%},
\]
где \(r\)~--- численное значение годовой ставки, записанное в процентах.

\blocktitle{Задача на закрепление}
Кредит \(4\) млн руб. выдан на \(8\) лет с равномерным погашением основного долга. Наибольший годовой платёж составляет не более \(0{,}9\) млн руб., а наименьший --- не менее \(0{,}54\) млн руб. Найдите \textbf{все возможные годовые процентные ставки}. Ответ запишите в процентах.
\end{document}